% ===================================================================
%  TABLEAU N° 13 — L'hôtel. --- Le restaurant. Le café.
%
%  Traduction, faite en 2026, en francais standard du Quebec, sur
%  l'ido de texto/io. Regles de la colonne : voir 10-tableau-01.tex.
%
%  LES NEUF CHOIX PROPRES A CE TABLEAU :
%    (12) bureau ................. réceptionniste
%    (4)  cabinet de lecture ..... salle de lecture
%    (26) carton à chapeau ....... boîte à chapeau
%    (38) appareil photographique . appareil photo
%    (55) bocks .................. verres de bière
%    (71) lorgnette .............. jumelles
%    (41) boîte à cirer .......... boîte à cirage
%    (20) bicyclette ............. vélo
%    (19) chasseur ............... groom
%
%  (19) ET (54) SONT LE MEME MOT EN IDO — « grumo » — et le fac-simile
%  les separe : « chasseur » pour celui que la caissiere envoie faire
%  une commission, « groom » pour celui qui porte les verres de biere.
%  On les reunit sous le second, comme aux tableaux 6 et 12 pour
%  robineto et trotuaro. Le premier mot, en 2026, ne nomme plus qu'un
%  homme au fusil.
%
%  LES REPAS DE CE TABLEAU SONT LA MEILLEURE ILLUSTRATION DE LA REGLE
%  POSEE AU TABLEAU 4, parce que l'ido y distingue lui-meme les deux :
%  « dejuneto » est le petit repas que la femme de chambre monte a la
%  chambre le matin — c'est le DEJEUNER ; « dejuno » est le repas de
%  la table d'hote, celui du menu, des tripes et du gigot — c'est le
%  DINER. Le fac-simile appelle les deux « dejeuner ».
%
%  « TABLE D'HOTE » RESTE, et il faut le dire parce que le mot a
%  change de metier des deux cotes. En 1926 c'etait la table commune
%  ou tous les pensionnaires mangeaient le meme repas a heure fixe ;
%  au Quebec de 2026 c'est le menu a prix fixe qu'un restaurant
%  affiche le midi. Les deux sens tiennent dans la phrase — un repas
%  servi a heure dite, sans choix — et le mot est vivant ici.
% ===================================================================

%%K t13-noto-1 noto
\VUnotes{143.2mm}{%
\textit{\VUgras{(*) Pour le détail de la chambre, voir le tableau
n\textsuperscript{o} 6.}}%
}

%%K t13-apar-1 apar
\VUsaut{3.0mm}
\VUtitre{15.0pt}{60}{TABLEAU N\textsuperscript{o} 13}
\VUsaut{1.6mm}
\VUpk{10.2pt}{40}{L'hôtel. --- Le restaurant.}
\VUsaut{2.0mm}
\VUpk{10.2pt}{40}{Le café.}
\VUsaut{3.4mm}

%%K t13-01-1 p
1. --- Arrivé hier soir à Royan, je suis descendu à l'« \textit{Hôtel de
l'Océan} », qu'un ami m'avait recommandé.

%%K t13-01-2 p
On m'a donné une jolie \VUgras{chambre} \textsuperscript{(2)}, au
\VUgras{deuxième étage} \textsuperscript{(1)}, où j'ai très bien dormi (*).

%%K t13-02-1 p
2. --- Ce matin, en sortant de ma chambre, où la \VUgras{femme de chambre}
\textsuperscript{(3)} avait apporté le déjeuner, je suis allé au
\VUgras{fumoir} \textsuperscript{(5)}, qui sert aussi de \VUgras{salle de
lecture} \textsuperscript{(4)}, pour parcourir les \VUgras{journaux}
\textsuperscript{(6)} en fumant une \VUgras{cigarette} \textsuperscript{(8)}.
Après ma lecture, je suis descendu par l'\VUgras{ascenseur}
\textsuperscript{(9)} et je suis allé me promener en ville.

%%K t13-03-1 p
3. --- Comme j'avais oublié de demander au \VUgras{réceptionniste}
\textsuperscript{(12)} de l'hôtel à quelle heure on servait les repas à la
\VUgras{table d'hôte}\textsuperscript{(11)}, je suis rentré tôt et
j'en ai profité pour aller prendre un bain dans la \VUgras{salle de bain}
\textsuperscript{(10)} de l'hôtel. Je suis ensuite revenu à la \VUgras{caisse}
\textsuperscript{(15)} pour téléphoner à un de mes amis. Mais le
\VUgras{téléphone} \textsuperscript{(13)} était occupé; voyant mon embarras, la
\VUgras{caissière} \textsuperscript{(14)}, qui inscrivait une \VUgras{facture}
\textsuperscript{(17)} sur le \VUgras{registre des voyageurs}
\textsuperscript{(16)}, a prié le \VUgras{concierge} \textsuperscript{(18)}
d'envoyer le \VUgras{groom} \textsuperscript{(19)} faire ma commission à
\VUgras{vélo} \textsuperscript{(20)}.

%%K t13-04-1 p
4. --- Je l'ai remerciée et je suis sorti pour donner un pourboire au
\VUgras{porteur} \textsuperscript{(24)} (le manœuvre) qui avait apporté de la
gare, sur sa \VUgras{voiture à bras} \textsuperscript{(25)}, ma \VUgras{boîte à
chapeau} \textsuperscript{(26)}, \VUgras{ma valise} \textsuperscript{(27)} et
\VUgras{mon sac de voyage} \textsuperscript{(28)}. Le \VUgras{facteur}
\textsuperscript{(21)}, qui venait d'arriver, m'a remis une \VUgras{lettre}
\textsuperscript{(22)}.

%%K t13-05-1 p
5. --- Je suis resté encore un moment sur le seuil de la porte, près de la
\VUgras{boîte aux lettres} \textsuperscript{(23)}, à regarder le mouvement de la
rue. Le \VUgras{conducteur} \textsuperscript{(30)} de l'\VUgras{omnibus}
\textsuperscript{(29)} chargeait sur sa voiture la \VUgras{malle}
\textsuperscript{(31)} d'un \VUgras{voyageur} \textsuperscript{(32)} auquel
l'\VUgras{hôtelier} \textsuperscript{(33)} faisait ses adieux. Un jeune
\VUgras{télégraphiste} \textsuperscript{(35)} apportait sans se presser une
\VUgras{dépêche} \textsuperscript{(34)} pour un client de l'hôtel; un
\VUgras{touriste} \textsuperscript{(36)}, son \VUgras{appareil photo}
\textsuperscript{(38)} en bandoulière, consultait son \VUgras{guide}
\textsuperscript{(37)}.

%%K t13-tit-1 sub
\VUsaut{4.0mm}
\VUpk{10.2pt}{40}{L'heure du dîner.}
\VUsaut{3.0mm}

%%K t13-06-1 p
6. --- Devant la grille, un placide \VUgras{commissionnaire}
\textsuperscript{(39)} sommeillait entre ses \VUgras{crochets}
\textsuperscript{(40)} et sa \VUgras{boîte à cirage} \textsuperscript{(41)},
pendant qu'une jeune \VUgras{blanchisseuse} \textsuperscript{(42)}, qui portait
un grand \VUgras{panier} \textsuperscript{(43)} de linge, riait de l'air
solennel du \VUgras{maître d'hôtel} \textsuperscript{(44)} qui se tenait près de
la porte.

%%K t13-07-1 p
7. --- La vue du \VUgras{cuisinier} \textsuperscript{(45)}, qui était sorti de
sa \VUgras{cuisine} \textsuperscript{(47)}, située au \VUgras{sous-sol}
\textsuperscript{(48)}, pour envoyer un \VUgras{marmiton}
\textsuperscript{(46)} faire une commission, m'a rappelé que l'heure du dîner
approchait, et je suis entré dans la salle à manger en traversant la cour, dans
laquelle un \VUgras{chauffeur} \textsuperscript{(50)} garait son
\VUgras{automobile} \textsuperscript{(49)}, tandis qu'un \VUgras{mécanicien}
\textsuperscript{(52)} réparait, suivant les indications du \VUgras{cycliste}
\textsuperscript{(53)}, un \VUgras{tricycle à pétrole} \textsuperscript{(51)}
qui venait d'avoir un accident.

%%K t13-08-1 p
8. --- J'ai demandé à un petit \VUgras{groom} \textsuperscript{(54)} qui portait
des \VUgras{verres de bière} \textsuperscript{(55)} si la table d'hôte était
sous la véranda et, sur sa réponse affirmative, j'ai pris place à une des tables
et je me suis fait servir un excellent repas.

%%K t13-noto-2 noto
\VUnotes{143.2mm}{%
\textit{\VUgras{(*) À l'aide de la liste des plats donnée dans le
vocabulaire, le professeur pourra facilement varier le menu ci-dessous.}}%
}

%%K t13-tit-2 sub
\VUsaut{4.0mm}
\VUpk{10.2pt}{40}{Un excellent dîner.}
\VUsaut{3.0mm}

%%K t13-09-1 p
9. --- Le garçon m'a apporté le menu (*) du dîner et la carte des vins. J'ai
choisi et j'ai commandé comme \VUgras{hors-d'œuvre} des \VUgras{crevettes} avec
du \VUgras{beurre} et, comme entrée, des \VUgras{tripes à la mode de Caen}. Je
n'ai pas pris de \VUgras{poisson}, mais j'ai demandé une portion de
\VUgras{gigot} d'agneau et, comme \VUgras{légume}, des \VUgras{épinards} à la
crème. Ensuite, comme aucun des \VUgras{entremets} ne me plaisait, je me suis
fait apporter divers desserts : du \VUgras{fromage}, des \VUgras{fruits} et des
\VUgras{biscuits}. J'y ai ajouté un excellent \VUgras{vin} de Saint-Émilion et,
très content de mon repas, j'ai quitté la table après avoir donné au
\VUgras{garçon} \textsuperscript{(57)} un beau \VUgras{pourboire}
\textsuperscript{(59)}.

%%K t13-10-1 p
10. --- Me voyant prêt à partir, le \VUgras{gérant} \textsuperscript{(60)} m'a
proposé de monter au balcon pour admirer le superbe panorama de l'\VUgras{Océan}
\textsuperscript{(62)}. On apercevait au loin de petites \VUgras{barques}
\textsuperscript{(63)} de pêche, des \VUgras{voiliers} \textsuperscript{(64)} et
un énorme \VUgras{paquebot} \textsuperscript{(65)} dont la silhouette noire se
détachait sur les \VUgras{nuages} \textsuperscript{(67)} blancs de
l'\VUgras{horizon} \textsuperscript{(66)}. Une brise légère faisait flotter le
\VUgras{drapeau} \textsuperscript{(70)} planté au-dessus de l'\VUgras{enseigne}
\textsuperscript{(69)} du \VUgras{hall} \textsuperscript{(68)}.

%%K t13-tit-3 sub
\VUsaut{3.0mm}
\VUpk{10.2pt}{40}{Les distractions.}
\VUsaut{3.0mm}

%%K t13-11-1 p
11. --- Le spectacle était si intéressant que j'ai décidé de monter demain sur
la terrasse du café en emportant des \VUgras{jumelles} \textsuperscript{(71)} ou
une \VUgras{longue-vue} \textsuperscript{(72)}.

%%K t13-12-1 p
12. --- En quittant le balcon, je me suis rendu au café de l'hôtel pour prendre
une \VUgras{consommation} \textsuperscript{(73)} en faisant une \VUgras{partie
de billard} \textsuperscript{(75)}. J'ai traversé sans m'arrêter la salle du
café, où beaucoup de consommateurs jouaient aux \VUgras{échecs}
\textsuperscript{(78)}, aux \VUgras{dames} \textsuperscript{(79)}, aux
\VUgras{dominos} \textsuperscript{(80)}, au \VUgras{trictrac}
\textsuperscript{(81)} ou aux \VUgras{cartes} \textsuperscript{(82)}, et je suis
monté directement à la \VUgras{salle de billard} \textsuperscript{(74)}.

%%K t13-13-1 p
13. --- Quand la partie a été finie, je suis descendu au rez-de-chaussée et j'ai
demandé au garçon une \VUgras{enveloppe} \textsuperscript{(84)}, du
\VUgras{papier} \textsuperscript{(83)}, un \VUgras{timbre-poste}
\textsuperscript{(85)}, une \VUgras{plume} \textsuperscript{(87)} et de
l'\VUgras{encre} \textsuperscript{(86)} pour écrire une lettre.

%%K t13-13-2 p
Puis j'ai appelé un \VUgras{cocher} \textsuperscript{(92)} dont la
\VUgras{voiture} \textsuperscript{(88)} s'était arrêtée devant le café. Je lui
ai demandé son prix à l'heure ou à la course et, quand nous nous sommes entendus
sur le prix, il m'a ouvert la \VUgras{portière} \textsuperscript{(89)} et m'a
installé. Il est remonté sur son \VUgras{siège} \textsuperscript{(91)}, a fait
partir les chevaux d'un léger coup de \VUgras{fouet} \textsuperscript{(93)} et
nous sommes partis pour une charmante promenade.
\VUsaut{6.0mm}
\VUfilet{22mm}
